% GENERATED FILE. Edit canonical lessons, then run npm run generate:books.
% canonical-source-hash: fnv1a64:b6be212adf9091d1
% canonical-lessons: SA-C537-pasakah, SA-C537-harayati, SA-C537-vancakah, SA-C537-lubdhah, SA-C537-dyutakarah

\chapter{A Game of Chance}
\label{ch:sa-a-game-of-chance}
\hlchaptermodality{537}

\begin{chapteropening}
\textbf{By the end of this chapter:} \emph{I can talk about dice, stakes and losing.}

Complete the last lesson of chapter 537: I can talk about dice, stakes and losing.
\end{chapteropening}

\section[\texorpdfstring{\sk{पाशकः} (pāśakaḥ)}{पाशकः (pāśakaḥ)}]{\sk{पाशकः} (pāśakaḥ) — a die}

\label{lesson:SA-C537-pasakah}



\begin{quote}
Before the new one: say the Sanskrit for by chance, then the Sanskrit for turns out well.
\end{quote}

\subsection*{You'll want to know: \sk{पाशकः}}
\textbf{\sk{पाशकः}} — \emph{pāśakaḥ} — \textquotedblleft{}a die\textquotedblright{}.

\begin{quote}
\sk{यदि} \sk{पाशकः} \sk{षट्} \sk{दर्शयति}, \sk{तर्हि} \sk{त्वं} \sk{जयसि।}
\end{quote}

\begin{quote}
\emph{yadi pāśakaḥ ṣaṭ darśayati, tarhi tvaṃ jayasi.}
\end{quote}

\begin{quote}
If the die shows six, you win.
\end{quote}

A masculine noun.

\begin{cousinweb}
\textbf{\sk{पाशकः}} (\emph{pāśakaḥ}) is built on \emph{pāśa}, \textquotedblleft{}a cord, a noose\textquotedblright{}, with the ending \emph{-ka}; MW gives \emph{pāśaka} as a snare, and also as a die (MW).
\end{cousinweb}

\begin{grammarlens}[title={What you've built}]
One more word for saying what happens if.
\end{grammarlens}

\subsection*{Your turn}
\begin{itemize}
\raggedright
  \item \emph{Say it:} \emph{pāśakaḥ}
  \item \emph{Say it:} \emph{pāśakaḥ}, once more
  \item \emph{Say it:} \emph{yadi pāśakaḥ ṣaṭ darśayati, tarhi tvaṃ jayasi.} — If the die shows six, you win.
  \item \emph{Recall:} say the Sanskrit for by chance, then the Sanskrit for turns out well, then say \emph{pāśakaḥ} again
\end{itemize}

\begin{tcolorbox}[breakable,colback=teal!4,colframe=teal!35!black,title={Before you move on}]
What does \sk{पाशकः} mean? (\textquotedblleft{}A die\textquotedblright{}.) Say it once more.
\end{tcolorbox}
\section[\texorpdfstring{\sk{हारयति} (hārayati)}{हारयति (hārayati)}]{\sk{हारयति} (hārayati) — loses at play}

\label{lesson:SA-C537-harayati}



\begin{quote}
Before the new one: say the Sanskrit for turns out well, then the Sanskrit for a die.
\end{quote}

\subsection*{You'll want to know: \sk{हारयति}}
\textbf{\sk{हारयति}} — \emph{hārayati} — \textquotedblleft{}loses at play\textquotedblright{}.

\begin{quote}
\sk{यदि} \sk{त्वं} \sk{द्यूतं} \sk{क्रीडसि}, \sk{तर्हि} \sk{त्वं} \sk{धनं} \sk{हारयसि।}
\end{quote}

\begin{quote}
\emph{yadi tvaṃ dyūtaṃ krīḍasi, tarhi tvaṃ dhanaṃ hārayasi.}
\end{quote}

\begin{quote}
If you gamble, you lose money.
\end{quote}

Past (he …): \sk{सः} \sk{सर्वं} \sk{धनं} \sk{हारितवान्}, \emph{saḥ sarvaṃ dhanaṃ hāritavān}.

Future: \sk{हारयिष्यति}, \emph{hārayiṣyati}.

\begin{cousinweb}
\textbf{\sk{हारयति}} (\emph{hārayati}) comes from the causative of the root √hṛ, \textquotedblleft{}to take, to carry off\textquotedblright{}: MW gives it as to have a thing taken from oneself, to lose it, above all at play (MW).
\end{cousinweb}

\begin{grammarlens}[title={What you've built}]
One more word for saying what happens if.
\end{grammarlens}

\subsection*{Your turn}
\begin{itemize}
\raggedright
  \item \emph{Say it:} \emph{hārayati}
  \item \emph{Say it:} \emph{hārayati}, once more
  \item \emph{Say it:} \emph{yadi tvaṃ dyūtaṃ krīḍasi, tarhi tvaṃ dhanaṃ hārayasi.} — If you gamble, you lose money.
  \item \emph{Recall:} say the Sanskrit for turns out well, then the Sanskrit for a die, then say \emph{hārayati} again
\end{itemize}

\begin{tcolorbox}[breakable,colback=teal!4,colframe=teal!35!black,title={Before you move on}]
What does \sk{हारयति} mean? (\textquotedblleft{}Loses at play\textquotedblright{}.) Say it once more.
\end{tcolorbox}
\section[\texorpdfstring{\sk{वञ्चकः} (vañcakaḥ)}{वञ्चकः (vañcakaḥ)}]{\sk{वञ्चकः} (vañcakaḥ) — a cheat}

\label{lesson:SA-C537-vancakah}



\begin{quote}
Before the new one: say the Sanskrit for a die, then the Sanskrit for loses at play.
\end{quote}

\subsection*{You'll want to know: \sk{वञ्चकः}}
\textbf{\sk{वञ्चकः}} — \emph{vañcakaḥ} — \textquotedblleft{}a cheat\textquotedblright{}.

\begin{quote}
\sk{यदि} \sk{सः} \sk{वञ्चकः} \sk{अस्ति}, \sk{तर्हि} \sk{अहं} \sk{न} \sk{क्रीडामि।}
\end{quote}

\begin{quote}
\emph{yadi saḥ vañcakaḥ asti, tarhi ahaṃ na krīḍāmi.}
\end{quote}

\begin{quote}
If he is a cheat, I won't play.
\end{quote}

A masculine noun.

\begin{cousinweb}
\textbf{\sk{वञ्चकः}} (\emph{vañcakaḥ}) comes from the causative of the root √vañc, \textquotedblleft{}to go crookedly\textquotedblright{}, which means \textquotedblleft{}to deceive\textquotedblright{}, with the agent ending \emph{-aka}; MW gives \emph{vañcaka} as deceiving, a deceiver, a cheat (MW).
\end{cousinweb}

\begin{grammarlens}[title={What you've built}]
One more word for saying what happens if.
\end{grammarlens}

\subsection*{Your turn}
\begin{itemize}
\raggedright
  \item \emph{Say it:} \emph{vañcakaḥ}
  \item \emph{Say it:} \emph{vañcakaḥ}, once more
  \item \emph{Say it:} \emph{yadi saḥ vañcakaḥ asti, tarhi ahaṃ na krīḍāmi.} — If he is a cheat, I won't play.
  \item \emph{Recall:} say the Sanskrit for a die, then the Sanskrit for loses at play, then say \emph{vañcakaḥ} again
\end{itemize}

\begin{tcolorbox}[breakable,colback=teal!4,colframe=teal!35!black,title={Before you move on}]
What does \sk{वञ्चकः} mean? (\textquotedblleft{}A cheat\textquotedblright{}.) Say it once more.
\end{tcolorbox}
\section[\texorpdfstring{\sk{लुब्धः} (lubdhaḥ)}{लुब्धः (lubdhaḥ)}]{\sk{लुब्धः} (lubdhaḥ) — covetous; greedy for gain}

\label{lesson:SA-C537-lubdhah}



\begin{quote}
Before the new one: say the Sanskrit for loses at play, then the Sanskrit for a cheat.
\end{quote}

\subsection*{You'll want to know: \sk{लुब्धः}}
\textbf{\sk{लुब्धः}} — \emph{lubdhaḥ} — \textquotedblleft{}covetous; greedy for gain\textquotedblright{}.

\begin{quote}
\sk{यदि} \sk{त्वं} \sk{लुब्धः} \sk{असि}, \sk{तर्हि} \sk{त्वं} \sk{सर्वं} \sk{धनं} \sk{हारयसि।}
\end{quote}

\begin{quote}
\emph{yadi tvaṃ lubdhaḥ asi, tarhi tvaṃ sarvaṃ dhanaṃ hārayasi.}
\end{quote}

\begin{quote}
If you are greedy, you lose all your money.
\end{quote}

\begin{cousinweb}
\textbf{\sk{लुब्धः}} (\emph{lubdhaḥ}) comes from the root √lubh, \textquotedblleft{}to go astray, to long for eagerly\textquotedblright{}, as its past participle; MW gives \emph{lubdha} as bewildered, and as greedy, covetous (MW).
\end{cousinweb}

\begin{grammarlens}[title={What you've built}]
One more word for saying what happens if.
\end{grammarlens}

\subsection*{Your turn}
\begin{itemize}
\raggedright
  \item \emph{Say it:} \emph{lubdhaḥ}
  \item \emph{Say it:} \emph{lubdhaḥ}, once more
  \item \emph{Say it:} \emph{yadi tvaṃ lubdhaḥ asi, tarhi tvaṃ sarvaṃ dhanaṃ hārayasi.} — If you are greedy, you lose all your money.
  \item \emph{Recall:} say the Sanskrit for loses at play, then the Sanskrit for a cheat, then say \emph{lubdhaḥ} again
\end{itemize}

\begin{tcolorbox}[breakable,colback=teal!4,colframe=teal!35!black,title={Before you move on}]
What does \sk{लुब्धः} mean? (\textquotedblleft{}Covetous; greedy for gain\textquotedblright{}.) Say it once more.
\end{tcolorbox}
\section[\texorpdfstring{\sk{द्यूतकारः} (dyūtakāraḥ)}{द्यूतकारः (dyūtakāraḥ)}]{\sk{द्यूतकारः} (dyūtakāraḥ) — a gambler}

\label{lesson:SA-C537-dyutakarah}



\begin{quote}
Before the new one: say the Sanskrit for a cheat, then the Sanskrit for covetous.
\end{quote}

\subsection*{You'll want to know: \sk{द्यूतकारः}}
\textbf{\sk{द्यूतकारः}} — \emph{dyūtakāraḥ} — \textquotedblleft{}a gambler\textquotedblright{}.

\begin{quote}
\sk{यदि} \sk{द्यूतकारः} \sk{सदा} \sk{जयति}, \sk{तर्हि} \sk{सः} \sk{वञ्चकः} \sk{भवेत्।}
\end{quote}

\begin{quote}
\emph{yadi dyūtakāraḥ sadā jayati, tarhi saḥ vañcakaḥ bhavet.}
\end{quote}

\begin{quote}
If a gambler wins every time, he may be a cheat.
\end{quote}

A masculine noun.

\begin{cousinweb}
\textbf{\sk{द्यूतकारः}} (\emph{dyūtakāraḥ}) is built from \emph{dyūta}, \textquotedblleft{}gambling, dice\textquotedblright{}, and \emph{kāra}, \textquotedblleft{}doing\textquotedblright{}, from the root √kṛ; MW gives \emph{dyūtakāra} as the same as \emph{dyūtakara}, a gambler (MW).
\end{cousinweb}

\begin{grammarlens}[title={What you've built}]
One more word for saying what happens if.
\end{grammarlens}

\subsection*{Your turn}
\begin{itemize}
\raggedright
  \item \emph{Say it:} \emph{dyūtakāraḥ}
  \item \emph{Say it:} \emph{dyūtakāraḥ}, once more
  \item \emph{Say it:} \emph{yadi dyūtakāraḥ sadā jayati, tarhi saḥ vañcakaḥ bhavet.} — If a gambler wins every time, he may be a cheat.
  \item \emph{Recall:} say the Sanskrit for a cheat, then the Sanskrit for covetous, then say \emph{dyūtakāraḥ} again
\end{itemize}

\begin{tcolorbox}[breakable,colback=teal!4,colframe=teal!35!black,title={Before you move on}]
What does \sk{द्यूतकारः} mean? (\textquotedblleft{}A gambler\textquotedblright{}.) Say it once more.
\end{tcolorbox}
